\documentclass[conference]{IEEEtran}
\usepackage{amsmath}
\usepackage{balance}
\begin{document}
\title{Dimensionality Reduction for Satellite Ship Classification}
\author{\IEEEauthorblockN{Taher Akolawala}}
\maketitle
\begin{abstract}
We evaluated whether dimensionality reduction improves classical ship classification on small satellite-image crops. Using 4,000 labeled, 80$\times$80 RGB images, we compared logistic regression and support vector machines (SVMs) on raw flattened pixels and representations produced by principal component analysis (PCA), Isomap, and locally linear embedding (LLE). On an 800-image holdout, PCA plus SVM achieved an F1 score of 0.8495, versus 0.7677 for the raw-pixel SVM; both LLE pipelines obtained zero F1. The selected PCA--SVM pipeline uses 150 components instead of 19,200 pixel features. These exploratory results come from one random image-level split and suggest that PCA can improve this particular classifier while reducing input dimensionality. Broader maritime use requires scene-aware validation.
\end{abstract}
\begin{IEEEkeywords}
satellite imagery, ship classification, principal component analysis, support vector machine, manifold learning
\end{IEEEkeywords}

\section{Introduction}
Ship recognition in overhead imagery is a useful binary classification problem: a small vessel can occupy only part of a crop, while water, shoreline, wakes, and lighting change the surrounding appearance. The public Ships in Satellite Imagery dataset contains 4,000 labeled crops from the San Francisco and San Pedro Bay areas, each with 80$\times$80 RGB pixels \cite{dataset}. A direct classical classifier receives 19,200 features per crop, large relative to the number of examples. This motivates a comparison of lower-dimensional representations before classification.

We asked whether reduced representations could preserve ship/no-ship discrimination while changing fit and prediction cost. To answer this, we trained comparable classifier pipelines on raw pixels and three reduction methods, selected settings by cross-validation, and evaluated the selected pipelines on a held-out test split. The study is exploratory and does not establish operational maritime performance.

\section{Data and Experimental Design}
The dataset arrays have shape $(4000,80,80,3)$ and contain 3,000 no-ship and 1,000 ship labels. We flattened each crop to a 19,200-dimensional vector and used an 80/20 random split with seed 91, yielding 3,200 training and 800 test images. The split was not stratified. We assembled preprocessing and estimators in scikit-learn pipelines, fitting scaling and dimensionality reduction within each training fold during cross-validation \cite{sklearn}.

The baseline estimators are logistic regression and an SVM. The reduced pipelines insert PCA, Isomap, or LLE before either classifier. Grid search optimizes F1 with three-fold cross-validation on the training data. Candidate PCA dimensions are 50, 80, 109, and 150; the saved search selects 150 for both PCA pipelines. A separate exploratory variance analysis identifies 109 components as approximately sufficient for 90\% explained variance. These numbers answer different questions: 109 describes a variance threshold, whereas 150 is the classifier dimension selected by the search.

For binary predictions, precision $P$ and recall $R$ are combined as
\begin{equation}
F_1=\frac{2PR}{P+R}.
\end{equation}
F1 is more informative than accuracy alone for the 1:3 ship-to-no-ship imbalance, though it still depends on the classifier's decision threshold. We did not evaluate transfer to a separate geographic region or use a scene-level holdout.

\section{Results}
\begin{table}[ht]
\caption{Results on the 800-image holdout}
\centering
\begin{tabular}{lcc}
\hline
Pipeline & F1 & Prediction time (s) \\
\hline
Raw SVM & 0.7677 & 3.6666 \\
Raw logistic regression & 0.7990 & 0.0566 \\
PCA + SVM & \textbf{0.8495} & 0.1613 \\
PCA + logistic regression & 0.8333 & 0.1485 \\
Isomap + SVM & 0.8324 & 1.7153 \\
Isomap + logistic regression & 0.8404 & 1.6647 \\
LLE + SVM & 0.0000 & 2.0357 \\
LLE + logistic regression & 0.0000 & 1.8648 \\
\hline
\end{tabular}
\end{table}
Table I shows that PCA plus SVM achieved the highest holdout F1 among the tested pipelines. The raw SVM took longer in this recorded prediction pass, although elapsed times depend on hardware and run-to-run variation and should not be treated as a stable latency benchmark. During three-fold cross-validation, PCA plus SVM reached a best F1 of 0.8447. Mean candidate fit times were 8.95 s in the PCA--SVM search and 17.95 s in the raw-SVM search. These values characterize the searches, rather than training time for a deployed selected model.

Both LLE pipelines obtained zero holdout F1. Their training-split confusion matrices show all-negative predictions, consistent with failure to recover the ship class at the chosen settings. Isomap performed better than LLE but did not exceed PCA plus SVM on the holdout. We do not infer test-set confusion counts from the training-split matrices. The holdout F1 values provide the clearest comparison among the evaluated pipelines.

\balance
\section{Limitations and Reproducibility}
This study is limited to one random image-level split. Crops from related geographic scenes may be correlated across training and test sets, so a scene-held-out evaluation would better probe transfer to new imagery. The split was not stratified, and we did not estimate confidence intervals or repeat the evaluation across random seeds. These constraints limit the precision and generality of the observed ranking.

Further work should report confusion matrices and precision--recall curves on the holdout, repeat the split across seeds, compare scene-level generalization, and rerun timing on a specified machine. A calibrated operating threshold would matter more than a single F1 score for any monitoring use. PCA and manifold methods should also be compared with a compact convolutional baseline before drawing conclusions about practical ship detection.

\section{Conclusion}
On the 800-image test split, PCA plus SVM attained 0.8495 F1 and exceeded the raw-pixel SVM's 0.7677. Its 150-component representation was substantially smaller than the original feature vector, while LLE failed to recover the ship class at the tested settings. PCA was the strongest pipeline in this exploratory comparison; broader deployment claims require repeated, scene-aware evaluation.

\begin{thebibliography}{2}
\bibitem{dataset} R. Hammell, ``Ships in Satellite Imagery,'' Kaggle dataset. [Online]. Available: \texttt{kaggle.com/datasets/rhammell/ships-in-satellite-imagery}
\bibitem{sklearn} F. Pedregosa \emph{et al.}, ``Scikit-learn: Machine learning in Python,'' \emph{J. Mach. Learn. Res.}, vol. 12, pp. 2825--2830, 2011.
\end{thebibliography}
\end{document}
